\ifdefined\gkver\else\def\gkver{student}\fi
\documentclass[\gkver]{gaokaozhenti}
\begin{document}

\chapter{2008年北京理综}
%% paperType: 理综物理部分

\begin{enumerate}
\item
%% number: 13
%% paperName: 2008年北京理综第 13 题 4 分
%% typeId: 1
%% type: 单选题
%% chapter: 光学
%% point: 光的干涉
%% method: 无
%% score: 4
%% degree: 850
%% duplicateId: 0
%% body:
下列说法中正确的是\xzanswer{A}

\fourchoices[answer=A]
{用分光镜观察光谱是利用光折射时的色散现象}
{用 X 光机透视人体是利用光电效应}
{光导纤维传输信号是利用光的干涉现象}
{门镜可以扩大视野是利用光的衍射现象}

%% answer: A
%% memo:
\memoanswer{本题原卷及材料中未提供详解，待补充。}

\item
%% number: 14
%% paperName: 2008年北京理综第 14 题 4 分
%% typeId: 1
%% type: 单选题
%% chapter: 原子和原子核
%% point: 核聚变
%% method: 无
%% score: 4
%% degree: 650
%% duplicateId: 0
%% body:
一个质子与一个中子聚变结合成一个氘核，同时辐射出一个 $\gamma$ 光子。已知质子、中子和氘核的质量分别为 $m_{1}$、$m_{2}$ 和 $m_{3}$，普朗克常量为 $h$，真空中的光速为 $c$，下列说法中正确的是\xzanswer{B}

\fourchoices[answer=B]
{核反应方程是 $\ce{^1_1 H + ^1_0 n -> ^3_1 H + \gamma}$}
{聚变反应中的质量亏损 $\Delta m=m_{1}+m_{2}-m_{3}$}
{辐射出的 $\gamma$ 光子的能量 $E=(m_{1}+m_{2}-m_{3})c$}
{$\gamma$ 光子的波长 $\lambda=\frac{h}{(m_{1}+m_{2}-m_{3})c^{2}}$}

%% answer: B
%% memo:
\memoanswer{本题原卷及材料中未提供详解，待补充。}

\item
%% number: 15
%% paperName: 2008年北京理综第 15 题 4 分
%% typeId: 1
%% type: 单选题
%% chapter: 气体性质
%% point: 阿伏伽德罗常数
%% method: 无
%% score: 4
%% degree: 650
%% duplicateId: 0
%% body:
假如全世界 60 亿人同时数 $1\Ug$ 水的分子数，每人每小时可以数 5000 个，不间断地数，则完成任务所需时间最接近的是（阿伏伽德罗常数 $N_{A}$ 取 $6\times10^{23}\Umoln$）\xzanswer{C}

\fourchoices[answer=C]
{10 年}
{1 千年}
{10 万年}
{1 千万年}

%% answer: C
%% memo:
\memoanswer{本题原卷及材料中未提供详解，待补充。}

\item
%% number: 16
%% paperName: 2008年北京理综第 16 题 4 分
%% typeId: 1
%% type: 单选题
%% chapter: 机械波
%% point: 波长频率波速
%% method: 无
%% score: 4
%% degree: 850
%% duplicateId: 0
%% body:
在介质中有一沿水平方向传播的简谐横波。一质点由平衡位置竖直向上运动，经 $0.1\Us$ 到达最大位移处，在这段时间内波传播了 $0.5\Um$，则这列波\xzanswer{D}

\fourchoices[answer=D]
{周期是 $0.2\Us$}
{波长是 $0.5\Um$}
{波速是 $2\Ums$}
{经 $1.6\Us$ 传播了 $8\Um$}

%% answer: D
%% memo:
\memoanswer{本题原卷及材料中未提供详解，待补充。}

\item
%% number: 17
%% paperName: 2008年北京理综第 17 题 4 分
%% typeId: 1
%% type: 单选题
%% chapter: 曲线运动
%% point: 人造地球卫星
%% method: 无
%% score: 4
%% degree: 650
%% duplicateId: 0
%% body:
据媒体报道，嫦娥一号卫星环月工作轨道为圆轨道，轨道高度 $200\Ukm$，运行周期 $127\Umin$。若还知道引力常量和月球平均半径，仅利用上述条件不能求出的是\xzanswer{B}

\fourchoices[answer=B]
{月球表面的重力加速度}
{月球对卫星的吸引力}
{卫星绕月运行的速度}
{卫星绕月运行的加速度}

%% answer: B
%% memo:
\memoanswer{本题原卷及材料中未提供详解，待补充。}

\item
%% number: 18
%% paperName: 2008年北京理综第 18 题 4 分
%% typeId: 1
%% type: 单选题
%% chapter: 交流电
%% point: 变压器
%% method: 无
%% score: 4
%% degree: 650
%% duplicateId: 0
%% body:
一理想变压器原、副线圈匝数比 $n_{1}:n_{2}=11:5$，原线圈与正弦交流电源连接，输入电压 $u$ 如图所示，副线圈仅接入一个 $10\UO$ 的电阻，则\xzanswer{D}

\onepicture[w=5.0cm]{figs/18.png}

\fourchoices[answer=D]
{流过电阻的电流是 $0.2\UA$}
{与电阻并联的电压表示数是 $100\sqrt{2}\UV$}
{经过 $1\Umin$ 电阻发出的热量是 $6\times10^{3}\UJ$}
{变压器的输入功率是 $1\times10^{3}\UW$}

%% answer: D
%% memo:
\memoanswer{本题原卷及材料中未提供详解，待补充。}

\item
%% number: 19
%% paperName: 2008年北京理综第 19 题 4 分
%% typeId: 1
%% type: 单选题
%% chapter: 磁场
%% point: 带电粒子在复合场中的运动
%% method: 无
%% score: 4
%% degree: 850
%% duplicateId: 0
%% body:
在如图所示的空间中，存在场强为 $E$ 的匀强电场，同时存在沿 $x$ 轴负方向、磁感应强度为 $B$ 的匀强磁场。一质子（电荷量为 $e$）在该空间恰沿 $y$ 轴正方向以速度 $v$ 匀速运动。据此可以判断出\xzanswer{C}

\onepicture[w=3.2cm]{figs/19.png}

\fourchoices[answer=C]
{质子所受电场力大小为 $eE$，运动中电势能减小，沿 $z$ 轴正方向电势升高}
{质子所受电场力大小为 $eE$，运动中电势能增大，沿 $z$ 轴正方向电势降低}
{质子所受电场力大小为 $evB$，运动中电势能不变，沿 $z$ 轴正方向电势升高}
{质子所受电场力大小为 $evB$，运动中电势能不变，沿 $z$ 轴正方向电势降低}

%% answer: C
%% memo:
\memoanswer{本题原卷及材料中未提供详解，待补充。}

\item
%% number: 20
%% paperName: 2008年北京理综第 20 题 4 分
%% typeId: 1
%% type: 单选题
%% chapter: 运动和力
%% point: 牛顿第二定律
%% method: 无
%% score: 4
%% degree: 450
%% duplicateId: 0
%% body:
有一些问题你可能不会求解，但你仍有可能对这些问题的解是否合理进行分析和判断。例如从解的物理量单位，解随某些已知量变化的趋势，解在一些特殊条件下的结果等方面进行分析，并与预期结果、实验结论等进行比较，从而判断解的合理性或正确性。

举例如下：如图所示，质量为 $M$、倾角为 $\theta$ 的滑块 A 放在水平地面上，把质量为 $m$ 的滑块 B 放在 A 的斜面上，忽略一切摩擦，有人求得 B 相对地面的加速度 $a=\frac{M+m}{M+m\sin^{2}\theta}g\sin\theta$，式中 $g$ 为重力加速度。

对于上述解，某同学首先分析了等号右侧量的单位，没发现问题，他进一步利用特殊条件对该解做了四项分析和判断，所得结论都是“解可能是正确的”，但是其中有一项是错误的，请你指出该项\xzanswer{D}

\onepicture[w=6.2cm]{figs/20.png}

\fourchoices[answer=D]
{当 $\theta=0^{\circ}$ 时，该解给出 $a=0$，这符合常识，说明该解可能对的}
{当 $\theta=90^{\circ}$ 时，该解给出 $a=g$，这符合实验结论，说明该解可能对的}
{当 $M\gg m$ 时，该解给出 $a=g\sin\theta$，这符合预期的结果，说明该解可能对的}
{当 $m\gg M$ 时，该解给出 $a=g/\sin\theta$，这符合预期的结果，说明该解可能对的}

%% answer: D
%% memo:
\memoanswer{本题原卷及材料中未提供详解，待补充。}

\item
%% number: 21
%% paperName: 2008年北京理综第 21 题 12 分
%% typeId: 4
%% type: 实验题
%% chapter: 力物体的平衡
%% point: 胡克定律
%% method: 无
%% score: 12
%% degree: 550
%% duplicateId: 0
%% body:
\begin{enumerate}
	\item
	用示波器观察某交流信号时，在显示屏上显示出一个完整的波形，如图。经下列四组操作之一，使该信号显示出两个完整波形，且波形幅度增大。此组操作是 \tkanswer[1.5cm]{C}（填选项代号）。
	
	\onepicture[w=4.0cm]{figs/21a.jpg}
	
	\fourchoices[answer=C]
	{调整 X 轴增益旋钮和竖直位移旋钮}
	{调整 X 轴增益旋钮和扫描微调旋钮}
	{调整扫描微调旋钮和 Y 轴增益旋钮}
	{调整水平位移旋钮和 Y 轴增益旋钮}
	\item
	某同学和你一起探究弹力和弹簧伸长的关系，并测量弹簧的劲度系数。做法是先将弹簧的一端固定在铁架台上，然后将最小刻度是毫米的刻度尺竖直放在弹簧的一侧，并使弹簧另一端的指针恰好落在刻度尺上。当弹簧自然下垂时，指针指示的刻度数值记作 $L_{0}$；弹簧下端挂一个 $50\Ug$ 砝码时，指针指示的刻度数值记作 $L_{1}$；弹簧下端挂两个 $50\Ug$ 砝码时，指针指示的刻度数值记作 $L_{2}$；……；挂七个 $50\Ug$ 砝码时，指针指示的刻度数值记作 $L_{7}$。
	
	①下表记录的是该同学已测出的 6 个值，其中有两个数值在记录时有误，它们的代号是 \tkanswer[1cm]{} 和 \tkanswer[1cm]{}。
	
	测量记录表：
	
	\begin{table}[htbp]
	\centering
	\begin{tblr}{|c|c|c|c|c|c|c|c|c|}
	\hline
	代表符号 & $L_{0}$ & $L_{1}$ & $L_{2}$ & $L_{3}$ & $L_{4}$ & $L_{5}$ & $L_{6}$ & $L_{7}$ \\
	\hline
	刻度数值（$\Ucm$） & 1.70 & 3.40 & 5.10 &  & 8.60 & 10.3 & 12.1 &  \\
	\hline
	\end{tblr}
	\end{table}
	
	②实验中，$L_{3}$ 和 $L_{7}$ 两个值还没有记录，请你根据上图将这两个测量值填入表中。
	
	\onepicture[w=9.0cm]{figs/21b.jpg}
	
	③为了充分利用测量数据，该同学将所测得的数值按如下方法逐一求差。分别计算出三个差值：$d_{1}=L_{4}-L_{0}=6.90\Ucm$，$d_{2}=L_{5}-L_{1}=6.90\Ucm$，$d_{3}=L_{6}-L_{2}=7.00\Ucm$，请你给出第 4 个差值 $d_{4}=$ \tkanswer[2.5cm]{} $=$ \tkanswer[1.5cm]{} $\Ucm$。
	
	④根据以上差值，可以计算出每增加 $50\Ug$ 砝码的弹簧平均伸长量 $\Delta L$，$\Delta L$ 用 $d_{1}$、$d_{2}$、$d_{3}$、$d_{4}$ 表示的式子为 $\Delta L=$ \tkanswer[3cm]{}；代入数据解得 $\Delta L=$ \tkanswer[1.5cm]{} $\Ucm$。
	
	⑤计算弹簧的劲度系数 $k=$ \tkanswer[1.5cm]{} $\UNm$（$g$ 取 $9.8\Umsq$）。
\end{enumerate}

%% answer:
\jdanswer{
\begin{enumerate}
	\item
	C
	
	\item
	① $L_{5}$，$L_{6}$；② 6.85，14.05；③ $L_{7}-L_{3}$，7.20；④ $\frac{d_{1}+d_{2}+d_{3}+d_{4}}{4\times4}$，1.75；⑤ 28
\end{enumerate}
}

%% memo:
\memoanswer{本题原卷及材料中未提供详解，待补充。}

\item
%% number: 22
%% paperName: 2008年北京理综第 22 题 16 分
%% typeId: 6
%% type: 计算题
%% chapter: 电磁感应
%% point: 电磁感应受力分析
%% method: 无
%% score: 16
%% degree: 650
%% duplicateId: 0
%% body:
均匀导线制成的单匝正方形闭合线框 abcd，每边长为 $L$，总电阻为 $R$，总质量为 $m$。将其置于磁感应强度为 $B$ 的水平匀强磁场上方 $h$ 处。如图所示，线框由静止起自由下落，线框平面保持在竖直平面内，且 cd 边始终与水平的磁场边界面平行。当 cd 边刚进入磁场时，

\begin{enumerate}
	\item
	求线框中产生的感应电动势大小；
	\item
	求 cd 两点间电势差的大小；
	\item
	若此时线框的加速度刚好为零，求线框下落的高度 $h$ 所应满足的条件。
\end{enumerate}

\onepicture[w=4.5cm, align=right]{figs/22.png}

%% answer:
\jdanswer{
\begin{enumerate}
	\item
	$E=BL\sqrt{2gh}$
	
	\item
	$U=\frac{3}{4}BL\sqrt{2gh}$
	
	\item
	$h=\frac{m^{2}gR^{2}}{2B^{4}L^{4}}$
\end{enumerate}
}

%% memo:
\memoanswer{
（1）线框速度 $v=\sqrt{2gh}$，感应电动势 $E=BLv=BL\sqrt{2gh}$。

（2）电流 $I=\frac{E}{R}$，两点间电压 $U=\frac{3}{4}IR=\frac{3}{4}BL\sqrt{2gh}$。

（3）安培力 $F=BIL=\frac{B^{2}L^{2}}{R}\sqrt{2gh}=mg$，得 $h=\frac{m^{2}gR^{2}}{2B^{4}L^{4}}$。
}

\item
%% number: 23
%% paperName: 2008年北京理综第 23 题 18 分
%% typeId: 6
%% type: 计算题
%% chapter: 功和能
%% point: 动能定理
%% method: 建立模型
%% score: 18
%% degree: 450
%% duplicateId: 0
%% body:
风能将成为 21 世纪大规模开发的一种可再生清洁能源。风力发电机是将风能（气流的动能）转化为电能的装置，其主要部件包括风轮机、齿轮箱、发电机等，如图所示。

\begin{enumerate}
	\item
	利用总电阻 $R=10\UO$ 的线路向用户输送风力发电机产生的电能，输送功率 $P_{\text{总}}=300\UkW$，输电电压 $U=10\UkV$，求导线上损失的功率与输送功率之比；
	\item
	风轮机叶片旋转所扫过的面积为风力发电机可接收风能的面积，设空气密度为 $\rho$，气流速度为 $v$，风轮机叶片长度为 $r$，求单位时间内流向风轮机的风能 $P_{\text{风}}$，在风速和叶片数确定的情况下，要提高风轮机单位时间接收的风能，简述可采取的措施；
	\item
	已知风力发电机输出的电功率 $P$ 与 $P_{\text{风}}$ 成正比，某风力发电机在风速 $v_{1}=9\Ums$ 时，能够输出电功率 $P_{1}=540\UW$。我国某地区风速不低于 $v_{2}=6\Ums$ 的时间每年约 5000 小时，试估算这台风力发电机在该地区的最小年发电量是多少千瓦时？
\end{enumerate}

\onepicture[w=5.5cm, align=right]{figs/23.png}

%% answer:
\jdanswer{
\begin{enumerate}
	\item
	$0.03$
	
	\item
	$P_{m}=\frac{1}{2}\pi\rho v^{3}r^{2}$，增加风轮机叶片长度、安装调向装置保持风轮机正面迎风等。
	
	\item
	$W=8\times10^{5}\UkWh$
\end{enumerate}
}

%% memo:
\memoanswer{
（1）导线上损失功率 $P=I^{2}R=P_{0}^{2}R/U^{2}=9\times10^{3}\UW$，损失功率与输送功率之比 $P/P_{0}=0.03$。

（2）风垂直流向风轮机时，提供的风能功率最大。单位时间内垂直流向叶片旋转面积的气体质量为 $\rho vS$，$S=\pi r^{2}$，风能的最大功率可表示为 $P_{m}=\frac{1}{2}\rho vSv^{2}=\frac{1}{2}\pi\rho v^{3}r^{2}$。可采取的措施：增加风轮机叶片长度，安装调向装置保持风轮机正面迎风等。

（3）按题意，风力发电机的输出功率 $P_{2}=\left(\frac{v_{2}}{v_{1}}\right)^{3}P_{1}=\left(\frac{6}{9}\right)^{3}\times540\UkW=160\UkW$，最小年发电量约为 $W=P_{2}t=160\times5000\UkWh=8\times10^{5}\UkWh$。
}

\item
%% number: 24
%% paperName: 2008年北京理综第 24 题 20 分
%% typeId: 6
%% type: 计算题
%% chapter: 运动和力
%% point: 动量守恒定律
%% method: 无
%% score: 20
%% degree: 450
%% duplicateId: 0
%% body:
有两个完全相同的小滑块 A 和 B，A 沿光滑水平面以速度 $v_{0}$ 与静止在平面边缘 O 点的 B 发生正碰，碰撞中无机械能损失。碰后 B 运动的轨迹为 OD 曲线，如图所示。

\begin{enumerate}
	\item
	已知滑块质量为 $m$，碰撞时间为 $\Delta t$，求碰撞过程中 A 对 B 平均冲力的大小；
	\item
	为了研究物体从光滑抛物线轨道顶端无初速下滑的运动，制做一个与 B 平抛轨迹完全相同的光滑轨道，并将该轨道固定在与 OD 曲线重合的位置，让 A 沿该轨道无初速下滑（经分析 A 在下滑过程中不会脱离轨道），
	
	a．分析 A 沿轨道下滑到任意一点时的动量 $p_{A}$ 与 B 平抛经过该点时的动量 $p_{B}$ 的大小关系；
	
	b．在 OD 曲线上有一点 M，O 和 M 两点的连线与竖直方向的夹角为 $45^{\circ}$，求 A 通过 M 点时的水平分速度和竖直分速度。
\end{enumerate}

\onepicture[w=5.5cm, align=right]{figs/24.png}

%% answer:
\jdanswer{
\begin{enumerate}
	\item
	$F=\frac{mv_{0}}{\Delta t}$
	
	\item
	a．$p_{A}<p_{B}$；b．$v_{xA}=\frac{2\sqrt{5}}{5}v_{0}$，$v_{yA}=\frac{4\sqrt{5}}{5}v_{0}$。
\end{enumerate}
}

%% memo:
\memoanswer{
（1）由动量守恒 $mv_{A}+mv_{B}=mv_{0}$，由机械能守恒 $\frac{1}{2}mv_{A}^{2}+\frac{1}{2}mv_{B}^{2}=\frac{1}{2}mv_{0}^{2}$，解得 $v_{A}=0$，$v_{B}=v_{0}$。对 B 有 $F\Delta t=mv_{0}$，所以 $F=\frac{mv_{0}}{\Delta t}$。

（2）a．设该点的竖直高度为 $d$，对 A 有 $E_{kA}=mgd$，对 B 有 $E_{kB}=mgd+\frac{1}{2}mv_{0}^{2}$，而 $p=\sqrt{2mE_{k}}$，所以 $p_{A}<p_{B}$。

b．对 B 有 $y=\frac{1}{2}gt^{2}$，$x=v_{0}t$，$y=\frac{g}{2v_{0}^{2}}x^{2}$，在 M 点 $x=y$，所以 $y=\frac{2v_{0}^{2}}{g}$；因轨迹相同，所以在任意点它们的速度方向相同。对 B 有 $v_{xB}=v_{0}$，$v_{yB}=\sqrt{2gy}=2v_{0}$，$v_{B}=\sqrt{5}v_{0}$；对 A 有 $v_{A}=\sqrt{2gy}=2v_{0}$。所以 $v_{xA}=v_{xB}\frac{v_{A}}{v_{B}}=\frac{2\sqrt{5}}{5}v_{0}$，$v_{yA}=v_{yB}\frac{v_{A}}{v_{B}}=\frac{4\sqrt{5}}{5}v_{0}$。
}

\end{enumerate}

\end{document}
